\capitulo{4}{Técnicas y herramientas}

En este capítulo se describen las técnicas y herramientas utilizadas durante el
desarrollo del proyecto. Se abordan tanto las metodologías empleadas como las
herramientas de software que facilitaron la implementación y gestión del mismo.

En la Figura~\ref{fig:tech_stack.png} se presenta un resumen conceptual del
ecosistema tecnológico empleado, categorizado por su función principal dentro
de la arquitectura.

\imagen{tech_stack.png}{Esquema de las principales tecnologías y librerías utilizadas.}{0.9}

\section{Marco de trabajo}

El desarrollo de \nombreProyecto{} se llevó a cabo aplicando prácticas ágiles
basadas en \acrfull{scrum}~\cite{agile_alliance} de forma adaptada al contexto
de un proyecto académico individual. Este enfoque permitió una gestión flexible
y adaptativa del proyecto. El trabajo se organizó en \gls{sprint} de dos
semanas, donde se definieron objetivos claros y se realizaron revisiones
periódicas para evaluar el progreso. Frente a enfoques estrictamente
secuenciales, este enfoque resulta especialmente adecuado en proyectos de
carácter experimental, como es el caso del modelado de temas, donde los
resultados dependen en gran medida de la calidad del preprocesado y de la
configuración de los modelos.

Para la planificación de tareas y realizar seguimiento del proyecto se utilizó
la aplicación web Zube~\cite{zube}. Esta herramienta se integra \gls{kanban}.
Los gráficos burndown se utilizaron junto con el tutor en la reuniones de
revisión de los sprints.

Para seguir el flujo de trabajo de \acrshort{scrum} y facilitar el proceso de
desarrollo, en este caso el tutor adoptó los roles de product owner y
\acrshort{scrum} master. Como evidencia de seguimiento, el tablero de Zube del
proyecto está disponible en \url{https://zube.io/tfg_topics/tfg-topics/}. De
forma complementaria, la trazabilidad de tareas e incidencias también puede
consultarse en las \textit{issues} del repositorio de GitHub del proyecto
\url{https://github.com/zcc1001/tfg-topics/issues}.

\section{ Plataforma de Gestión de Ciclo de Vida del Software (DevOps)}

Para la gestión del proyecto \nombreProyecto{} se adoptó la plataforma integral
de GitHub. Esta infraestructura se configuró como núcleo operativo para
soportar las prácticas del proceso de desarrollo adoptado. Las evidencias de
ejecución pública se encuentran accesibles en el repositorio oficial del
proyecto ejecución pueden consultarse en el repositorio oficial del proyecto
\url{https://github.com/zcc1001/tfg-topics}.

A nivel operativo, el ecosistema de trabajo se articuló en las siguientes
dimensiones:

\begin{itemize}
    \item \textbf{Gestión de versiones de código:} Utilizando Git como sistema de control de versiones distribuido, se aplicó un modelo de ramificación que favoreció la trazabilidad de cambios, el aislamiento por funcionalidades y la estabilidad de la rama principal~\cite{GitReference}.
    \item \textbf{Gestión de tareas y planificación ágil:} El seguimiento del \textit{backlog} y de las iteraciones se instrumentó con GitHub Issues y su integración con los tableros Kanban de Zube, facilitando la vinculación entre tareas, incidencias y \textit{commits}.
    \item \textbf{Integración continua y control de calidad:} Mediante GitHub Actions  se automatizaron los flujos de trabajo textit{workflows} para ejecutar baterías de pruebas unitarias (\textit{testing}) y análisis estático de calidad de código en cada solicitud de extracción (\textit{pull request}), asegurando la calidad del software.
    \item \textbf{Construcción de releases y despliegue:} Se definieron canalizaciones automatizadas para empaquetar versiones estables y sostener el despliegue continuo de los componentes de la aplicación.
    \item \textbf{Gestión del conocimiento y documentación:} Se utilizó GitHub Wiki para centralizar documentación técnica y funcional del sistema, estructurando de manera exhaustiva tanto el manual del programador —detallando la arquitectura hexagonal y la API— como el manual de usuario para la explotación de la interfaz en Streamlit.
    \item \textbf{Entornos de desarrollo virtualizados (GitHub Codespaces):} Se configuraron \textit{devcontainers} para habilitar entornos estandarizados  e independientes de la máquina local mediante GitHub Codespaces. Esta aproximación de \acrfull{eac} eliminó los conflictos de configuración de dependencias de Python y las librerías de modelado de temas. Asimismo, facilitó la comunicación síncrona y asíncrona con el tutor del proyecto, permitiendo desplegar de manera inmediata entornos efímeros pero completamente funcionales para testear y validar de forma conjunta los nuevos incrementos de software antes de su consolidación en producción.
    \item \textbf{Ingeniería de software asistida por IA:} GitHub Copilot se empleó como soporte de asistencia técnica en tareas de codificación, refactorización y revisión crítica de código, manteniendo supervisión humana en todo momento. Como evidencia, la incidencia pública \#63 documenta su uso en revisiones y propuestas de mejora: \url{https://github.com/zcc1001/tfg-topics/issues/63}.
\end{itemize}

\section{Herramientas de desarrollo}

\subsection{Entorno de desarrollo}

El entorno de desarrollo utilizado fue \textbf{Visual Studio Code}, un editor
de código fuente ligero y gratuito desarrollado por
Microsoft~\cite{QueEsVisualStudioCode}. Este entorno ofrece una amplia gama de
extensiones y funcionalidades que facilitan la escritura, depuración y gestión
de código en múltiples lenguajes de programación. Fue una herramienta esencial
para el desarrollo de este proyecto gracias a su amplia integración con
\textit{plugins} que ayudan al desarrollo y depuración con el lenguaje de
programación Python. Además,Visual Studio Code cuenta con integración nativa
con GitHub, lo que facilita la gestión de versiones. Se valoró el uso de otros
entornos como \textbf{PyCharm}~\cite{PyCharmTodoSobre2022}, pero se optó por
Visual Studio por su ligereza y flexibilidad.

También se utilizó \textbf{Jupyter
    Notebooks}~\cite{ProjectJupyterDocumentation} para la experimentación y
análisis de datos, ya que permite combinar código, visualizaciones y
documentación en un solo documento interactivo.

\subsection{Lenguajes de programación}

El lenguaje de programación principal utilizado en el desarrollo del proyecto
fue \textbf{Python}~\cite{31019Documentation}. Python es un lenguaje de alto
nivel, interpretado y de propósito general, conocido por su sintaxis clara. Su
elección se basa en la amplia gama de bibliotecas y \gls{framework} para el
procesamiento de datos y modelado de temas disponibles.
\subsubsection{Bibliotecas y frameworks}

El núcleo del proyecto se apoya en técnicas de \Acrshort{nlp} aplicadas al
análisis de documentos. Para el procesamiento de datos y modelado de temas se
utilizaron diversas bibliotecas de Python, entre las que destacan:

\begin{itemize}
    \item \textbf{Pandas}~\cite{UserGuidePandas}: Biblioteca fundamental para la manipulación y análisis de datos.
    \item \textbf{NumPy}~\cite{NumPyUserGuide}: Biblioteca para computación científica y manejo de arreglos multidimensionales.
    \item \textbf{Gensim}~\cite{GensimTopicModelling}: Biblioteca especializada en modelado de temas y procesamiento de lenguaje natural.
    \item \textbf{NLTK}~\cite{NLTKNaturalLanguage}: Biblioteca para el procesamiento de lenguaje natural, usado para tareas básicas tokenización y \textit{stopwords}.
    \item \textbf{spaCy}~\cite{SpaCy101Everything}: Biblioteca avanzada para procesamiento más avanzado de lenguaje natural.
\end{itemize}

Sobre los textos preprocesados se aplican técnicas de modelado de temas. Se han
estudiado y utilizado distintos enfoques, desde modelos probabilísticos
clásicos hasta métodos más recientes basados en representaciones vectoriales.
La comparación entre estas técnicas permite analizar sus fortalezas y
limitaciones en función de criterios como la coherencia semántica de los temas,
la interpretabilidad y el coste computacional.

Se utilizaron las siguientes bibliotecas de Python para el modelado de temas:
\begin{itemize}
    \item \acrshort{lda}~\cite{LdaTopicModelingb}: Implementación del modelo Latent Dirichlet Allocation para descubrir temas en grandes colecciones de documentos.
    \item \textbf{BERTopic}~\cite{QuickStartBERTopic}: Biblioteca que utiliza modelos de lenguaje basados en transformadores para generar representaciones de documentos y realizar clustering.
    \item \textbf{Top2Vec}~\cite{angelovDdangelovTop2Vec2025a}: Biblioteca que combina incrustaciones de documentos y clustering para identificar temas en grandes conjuntos de datos.
    \item \textbf{FASTopic}~\cite{FastopicFASTopic}: Biblioteca que utiliza técnicas de reducción de dimensionalidad y clustering para el modelado eficiente de temas en grandes conjuntos de datos.
\end{itemize}

En cuanto a la experimentación con modelos y la optimización de
hiperparámetros, se ha utilizado
\textbf{Optuna}~\cite{OptunaHyperparameterOptimization} como herramienta
principal para la exploración del espacio de configuraciones. Esta biblioteca
permite definir procesos de búsqueda automatizados y eficientes y facilita la
comparación objetiva de los resultados mediante métricas cuantitativas como la
coherencia de los temas. Frente a la selección manual de parámetros, este
enfoque resulta más riguroso y contribuye a la reproducibilidad y trazabilidad
de los experimentos realizados.

Finalmente, para la presentación y exploración de los resultados se ha
desarrollado una interfaz web utilizando la librería
\textbf{Streamlit}~\cite{StreamlitDocsa}. Su elección se debe principalmente a
que esta herramienta permite construir interfaces web interactivas de forma
ágil sin entrar en detalles complejos de desarrollo web. Streamlit facilita la
creación de paneles donde se pueden cargar datos y ajustar parámetros
facilitando la visualización de los modelos generados y de sus métricas
asociadas.

\subsection{Otras herramientas}

Además de las herramientas mencionadas anteriormente, se utilizaron diversas
herramientas de apoyo orientadas a la gestión del entorno de ejecución y a la
calidad del software desarrollado.

En primer lugar, se ha utilizado \textbf{Docker}~\cite{DockerManuals} y
\textbf{Docker Compose} como tecnología de \gls{contenedorizacion}. Docker
permite encapsular la aplicación junto con todas sus dependencias en entornos
reproducibles y aislados, lo que simplifica tanto el despliegue como la
ejecución del sistema en diferentes máquinas. Este enfoque resulta
especialmente útil ya que reduce problemas derivados de incompatibilidades de
versiones y facilita la portabilidad y reproducibilidad de los experimentos.

Por otro lado, se han utilizado sistemas orientados a garantizar la calidad del
código. En primer lugar, para el formateo automático del código se ha utilizado
\textbf{Black}, lo que asegura un estilo de código uniforme en todo el
proyecto. La organización y limpieza de importaciones se ha gestionado usando
\textbf{isort}, mientras que el análisis estático y la detección de posibles
errores o incumplimientos de estilo se han realizado con \textbf{flake8}. Estas
herramientas mejoran la legibilidad del código y reducen la aparición de
errores durante el desarrollo. A modo de síntesis, en la
Tabla~\ref{tabla:flujo_desarrollo_validacion} se detalla el flujo integral de
desarrollo, control de versiones y validación continua aplicado.

\tablaSmall{Flujo de desarrollo, control de versiones y validación continua.}
{p{0.22\textwidth} p{0.38\textwidth} p{0.30\textwidth}}
{flujo_desarrollo_validacion}
{
    \textbf{Etapa} & \textbf{Acciones principales} & \textbf{Herramientas} \\
}
{
    Planificación y seguimiento &
    Definición de tareas, priorización y revisión periódica del avance. &
    Zube, GitHub Issues. \\
    Implementación & Desarrollo incremental por módulos y ejecución local o
    contenedorizada. & Python, Visual Studio Code, Docker, Makefile. \\
    Control de versiones & Versionado del código, integración de cambios y
    trazabilidad de decisiones. & Git, GitHub. \\
    Calidad estática & Formateo, orden de importaciones, análisis de estilo y
    verificación de tipos. & Black, isort, flake8, mypy, pre-commit, SonarQube. \\
    Pruebas & Ejecución de pruebas unitarias por módulo y validación funcional de
    cambios. & Pytest. \\
    Integración y despliegue & Ejecución reproducible del \gls{pipeline} completo en
    distintos modos. & Docker Compose, \acrshort{ghcr}, Makefile, GitHub actions. \\
    Documentación & Redacción de memoria técnica y mantenimiento de documentación
    del proyecto. & LaTeX, Visual Studio Code, Wiki, markdown. \\
}

\section{Documentación}

Para la creación de la documentación del proyecto se utilizó
\textbf{LaTeX}~\cite{LaTeXDocumentPreparation}. LaTeX es especialmente adecuado
para la creación de documentos técnicos y científicos debido a su capacidad
para manejar fórmulas matemáticas complejas, referencias bibliográficas y
estructuras avanzadas de documentos. Además, LaTeX permite separar el contenido
de la presentación, facilitando la gestión y el formato del documento final. Se
optó por LaTeX frente a procesadores de texto tradicionales como Microsoft Word
o Google Docs, ya que estos últimos pueden presentar limitaciones en la gestión
de referencias y fórmulas complejas. Además, la universidad proporciona
plantillas pre formateadas y estandarizadas para la creación de la memoria en
LaTeX. Como editor de LaTeX se utilizó Visual Studio Code con la extensión
\textbf{LaTeX Workshop}~\cite{yuJamesYuLaTeXWorkshop2025}.

